%The existence of new physics Beyond the Standard Model (BSM) is backup up by strong experimental evidence and several theoretical indications. In the last few years, the possibility of having light new degrees of freedom, coupled to the SM only very feebly, has drawn significant attention in the community. In this regard, the use of high-intensity beam-dump experiments allows to explore new regions of parameter space for new particles that are light enough so they can be produced in the decay of unstable mesons produced in the target. Very recently it was pointed out that the CERN ProtoDUNEs may be used as an opportunistic beam-dump experiment precisely for this purpose, given their location with respect to the  upstream SPS T2 secondary beam target in the North Area at CERN. 

High-intensity beam-dump experiments currently lead the limits on the existence of new feebly interacting particles with masses in the MeV-GeV regime, which appear in many well-motivated  extensions Beyond the Standard Model (BSM).   
%
In this note we take the first steps towards demonstrating if ProtoDUNE detectors exposed to the SPS secondary beams can act as a beam-dump facility. Before any search for new BSM can be attempted, it is imperative to demonstrate that standard weakly interacting particles (namely neutrinos) coming from the T2 target can be observed. This is specially challenging given the relatively long SPS beam spill (of approximately 4.8~s) and the drift time of the detector (on the order of several ms).
%The analysis presented in this technical note constitutes an initial step in this direction and serves as a proof of concept. 
The analysis presented in this technical note focuses on demonstrating two key requirements for this task: (i) that the detector can be triggered on beam-related events, and (ii) that cosmic-ray events can be efficiently removed with appropriate selection cuts. In addition, we report the first observation of potentially accelerator-produced neutrino events in ProtoDUNE. Specifically, our results use a small subset of data recorded parasitically at the conclusion of the 2024 ProtoDUNE Horizontal Drift detector campaign, obtained with a dedicated trigger and offline software filter. 

The analysis relies on a time-based selection exploiting the correlation between the SPS beam spill and reconstructed activity in the detector, combined with topological and directional criteria designed to suppress cosmic-ray–induced backgrounds. Events are also classified according to the instantaneous trigger primitive (TP) rates, which varied during data taking. The analysis is therefore performed separately for periods characterized by low-TP and high-TP rates. We have checked that the event selection is capable of removing the additional contribution in periods with high TP rates, rendering the two final samples compatible. 

A clear excess above background is observed for events pointing along the SPS/T2 beam direction, compatible with the expectation from neutrino interactions with energies roughly above 10~GeV. Using all data that pass our selection cuts (summarized in Tab.~\ref{tab:post_cut_numbers_all_tps}) its associated statistical significance stands at $6.35~\sigma$. As an additional check, we have also computed the significance  restricting the sample to those events recorded in periods with low TP rates. A clear excess is also observed in this case, albeit with lower significance ($3.85~\sigma$) which may be attributed to the smaller size of the sample. 

Overall, our results seem to indicate that a search for a BSM signal using the T2 target as a beam-dump is feasible, in spite of the challenges outlined above. As a by-product of our analysis, we have also obtained a sample of neutrino interactions in LAr, which can be used for further studies. However, further steps are needed to determine that these events stem from neutrinos produced in the T2 target area. In particular, it should be demonstrated that they do not originate from the activity along H2/H4 beamlines. As discussed in this document, several experiments occupied the secondary beamlines with different beam requirements. For that purpose, different settings are used in the magnets surrounding the target and along the H2 and H4 beamlines. Dedicated studies are ongoing to demonstrate if the neutrino-like events can arise from high-energy pions or kaons produced and transported in the line. A run at PD-VD motivated by these results is requested to unambiguously disentangle the origin of these events. 

Finally, it is also worth mentioning that the measured event rate remains lower than the expectation from Monte Carlo simulations. This discrepancy may be attributed to the limited size of the sample (subject to large statistical errors) as well as to systematic uncertainties related to the neutrino flux (including hadron production), interaction cross sections, beamline instrumentation, and detector response. A detailed evaluation of these effects will be required in future analyses, as well as for any subsequent BSM searches.
